% ======================================================================
% SECTION 2: METHODS
% Final prose. Documents AI generations, repository inputs, build, and
% CI compatibility. Discusses simulation type (site vs sponsor) for each
% of the four simulations.
% ======================================================================

\subsection{AI Generations and Author Roles}
\label{sec:methods-ai}

Primary AI generations for the original paper template and the
subsequent population of that template into the full prose paper were
performed by Claude Code Opus 4.7 Max with a 1M token context window
using the Max 5x plan \cite{claude-opus-47, claude-code}. The
template-creation pass produced the LaTeX skeleton, the style file,
the bibliography, the section instruction blocks, and the prose
Methods section as a single multi-commit pull request to the
\path{kevinkawchak/physical-ai-oncology-trials} repository
\cite{main-repo}. The current pass (full-paper population) replaces
each bracketed instruction block with senior-author-quality final prose
and ships an Overleaf-ready ZIP at
\path{new-trial/national-24-7-trial/paper/full-paper/}.

The current first prompt was written by the author of the paper.\
ChatGPT Thinking 5.5 with Extended Thinking and Deep Research, accessed
through ChatGPT Plus, was used to generate both the Background-A and
Background-B summaries that seed the bibliography and the Introduction
\cite{chatgpt-thinking}. Claude Sonnet 4.6 with Adaptive Thinking was
used to produce three chunks and a README file separately for both
Background-A and Background-B, yielding the eight Markdown source
files at \path{new-trial/national-24-7-trial/Background-A/} and
\path{new-trial/national-24-7-trial/Background-B/}
\cite{claude-sonnet-46}. Google Gemini AI Overview, accessed through
\path{google.com} search, was used for additional search assistance
during the source-gathering phase \cite{google-gemini-overview}. The
author performed the senior-author white-space and table-width
formatting pass after the body sections were populated, applying the
formatting rules captured at the top of \path{main.tex}.

\subsection{Simulation Type: Clinical Trial Sites and Sponsors}
\label{sec:methods-sim-type}

The four simulations referenced throughout this paper split between
clinical trial sites and clinical trial sponsors. Simulations 1 and 2
are site-side simulations: Simulation 1 is the Continuous National
24/7 RTCT across four sites (SITE-A Houston, SITE-B Philadelphia,
SITE-C Boston, SITE-D Texas Medical Center), and Simulation 2 is the
single-patient 10-stage journey of patient PAT-2026-0042 through one
site's protocol from pre-screening through closeout. Simulations 3
and 4 are sponsor-side simulations: Simulation 3 is the 24-hour
autonomous sponsor with 53 core agents, and Simulation 4 is the
168-hour 7-day sponsor extension with local verification on smaller agent scripts with a
constrained Core i5-6200U laptop running smaller agentic scripts. Both perspectives are needed for a
continuous trial. Sites generate the patient and robot signal streams,
and sponsors generate the governance, regulatory, and decision
streams; the FDA RTCT pilot ingests both as a single real-time view.

\subsection{Repository Inputs}
\label{sec:methods-inputs}

The four simulations referenced throughout this paper are located at
the following repository paths and are documented in the Results
section:

\begin{itemize}[leftmargin=*]
\item Simulation 1 (clinical trial site) - Continuous Real-Time
Clinical Trial: \path{new-trial/national-24-7-trial/hour-00/} through
\path{hour-55/} (56 hours included in the paper), with each hour
folder containing four Markdown files (\texttt{simulation.md},
\texttt{robot\_logs.md}, \texttt{patient\_records.md},
\texttt{psl\_scores.md}) and three ASCII diagrams
(\texttt{diagram\_facility.txt}, \texttt{diagram\_patient\_flow.txt},
\texttt{diagram\_robot\_status.txt}). See
\cite{kawchak_2026_19176370, kawchak_2026_19244918} for the related
site documentation and national platform packages, and
\cite{kawchak_2026_18810541} for the patient instructions package.

\item Note on \path{extra-hours/}: the 28 hour folders
\path{new-trial/national-24-7-trial/extra-hours/hour-56/} through
\path{hour-83/} contain only approximated data and approximated
diagrams due to extended AI run time during the original cloud
generation pass, and are therefore \emph{not included} in the
full-paper text or analyses.  

\clearpage

\item Simulation 2 (clinical trial site) - Single-Patient 10-Stage
Journey: \path{patient-journey/stage_01_prescreening.py} through
\path{patient-journey/stage_10_closeout.py}, with the prior paper at
\path{patient-journey/paper/patient_journey_paper.tex}. The
orchestrators are \path{patient-journey/master_journey.py} and
\path{patient-journey/patient_state.py} \cite{kawchak_2026_19119939}.

\item Simulation 3 (clinical trial sponsor) - 24-Hour Autonomous
Sponsor: \path{sponsor/final_paper/scripts/hourly/sponsor_hour_00.py}
through \path{sponsor_hour_23.py}, with 53 core agents under
\path{sponsor/final_paper/scripts/core_agents/} and 75 ASCII diagrams
under \path{sponsor/final_paper/scripts/diagrams/}, plus 24 JSON
outputs at \path{sponsor/final_paper/scripts/hourly/output/} and the
cumulative \path{sponsor_24h_summary.json}
\cite{kawchak_2026_19396256}.

\item Simulation 4 (clinical trial sponsor) - 168-Hour 7-Day
Extension: \path{sponsor/final_paper/168_hours/day_01/} through
\path{day_07/}, with hourly scripts at
\path{day_NN/hourly/sponsor_hour_NNN.py} (168 scripts total), 168
JSON outputs at \path{day_NN/hourly/output/}, and 525 ASCII diagrams
at \path{day_NN/diagrams/}. The local-verification run is documented
at
\path{sponsor/final_paper/168_hours/instructions/core_i5_6200u_4gb/}
\cite{kawchak_2026_19396256}.
\end{itemize}

The deep-research summaries that anchor the Introduction are located
at \path{new-trial/national-24-7-trial/Background-A/} (three Markdown
chunks plus a README, 17 BibTeX entries) and
\path{new-trial/national-24-7-trial/Background-B/} (three Markdown
chunks plus a README, 9 BibTeX entries; one entry, Yoo2025SCORPIO, is
shared with Background-A and is included once in the bibliography).
The FDA April 28, 2026 source release is located at
\path{new-trial/national-24-7-trial/FDA-April-2026/FDA_RealTime_Clinical_Trials.md}
\cite{fda2026realtime}.

\subsection{Build and Verification}
\label{sec:methods-build}

The LaTeX paper compiles in Overleaf and locally with the standard
four-pass sequence:
\begin{verbatim}
pdflatex main.tex
bibtex main
pdflatex main.tex
pdflatex main.tex
\end{verbatim}

The bibliography uses the \texttt{ieeetr} style. Every reference
includes a DOI where one exists and a clickable URL through the
\texttt{note} field.\ The full-paper directory ships an Overleaf-ready
ZIP at
\path{new-trial/national-24-7-trial/paper/full-paper/LaTeX_Source_Files.zip}
that contains \texttt{main.tex}, \texttt{new\_paper.sty},
\texttt{references.bib}, \texttt{orcid\_icon.png}, the
\texttt{sections/} directory, and \texttt{README.md}.

\subsection{Continuous Integration Compatibility}
\label{sec:methods-ci}

The full-paper directory contains only LaTeX, Markdown, ZIP, and PNG
assets. No new Python files are introduced by the full-paper pass, so
the repository's existing \texttt{ruff} configuration in
\path{ruff.toml} and the \texttt{lint-and-format} workflow at
\path{.github/workflows/ci.yml} across Python 3.10, 3.11, and 3.12
remain green. The lint-and-format workflow runs \texttt{ruff format
--check} and \texttt{ruff check} on the existing Python tree, and the
pull request that delivers this full paper does not change any Python
file. Should a future pass introduce a helper Python script for
diagram generation or for repository indexing, that script must be
formatted with \texttt{ruff format} and pass \texttt{ruff check} under
the repository's 120-character line length and \texttt{[E, F, W]} rule
selection.

\clearpage

\subsection{Reproducibility: Cloud Only Versus Cloud Plus Local}
\label{sec:methods-reproducibility}

The four simulations differ in where their compute lives, which in
turn drives different reproducibility advantages and disadvantages.

\textbf{Cloud-only simulations.} Simulations 1, 2, and 3 ran entirely
on cloud compute, with all outputs (Markdown, ASCII, Python source,
JSON) committed back to the GitHub repository. The advantage is that
the cloud has high throughput - Simulation 1 generated 56 hours
(included in this paper) of minute-resolution narrative in roughly 4
to 5 hours of wall-clock cloud time, far faster than a local
generation pass. The disadvantage is that re-running a cloud-only
simulation requires cloud compute access, which is not always
available to a clinical reader who wishes to verify the artifacts
independently.

\textbf{Cloud-plus-local simulation.} Simulation 4 paired the cloud
generation with a 168-hour local verification of smaller agent scripts run on a 2015
Intel Core i5-6200U laptop with 4 GB of RAM running Windows 10 Pro.\
The cloud generated the 168 Python scripts and the 525 ASCII diagrams,
and the local laptop re-executed those scripts to produce a
\path{sponsor_168h_summary.json} that reflects the cloud output on a smaller scale. The
advantage is high reproducibility: anyone with a 4 GB Windows laptop
can re-execute the simulation and confirm the cumulative metrics. The
disadvantage is that the local verification was partial. Not every
script executed end-to-end on the constrained hardware due to thermal
throttling and operating-system-specific quirks (Windows Update,
antivirus exclusion, pagefile sizing), and downstream analytics must
account for the partial-verification status.

\subsection{Code-Based Versus Text-Only Simulations}
\label{sec:methods-code-vs-text}

The four simulations also differ in whether their primary output is
Python or only Markdown and ASCII text.

\textbf{Text-only simulation (Simulation 1).} Simulation 1 emits only
Markdown narrative and ASCII diagrams; no Python execution surface is
generated. The benefit is auditability: a clinical reader can read
\path{hour_NN_simulation.md} and the three ASCII diagrams in any text
viewer, with no Python interpreter required. The cost is that
downstream automation cannot easily consume the output - every
downstream agent must parse the Markdown and ASCII, which is brittle.

\textbf{Code-based simulations (Simulations 2, 3, 4).} The
patient-journey simulation (10 Python stage scripts), the 24-hour
sponsor (24 hourly Python scripts plus 53 agent modules), and the
168-hour sponsor (168 hourly Python scripts) all emit Python code as
the primary artifact. The benefit is that downstream agents can
import the modules directly, which gives the simulation a structured
execution surface that text-only simulations lack. The cost is that
a clinical reader who wants to verify the output must run the Python
code, which raises the technical barrier to verification.

The two approaches are complementary.\ A complete continuous-trial
deployment would use text-only artifacts (Simulation 1's narrative)
to document the patient and robot timeline for clinical readers, and
code artifacts (Simulations 2 to 4) to drive downstream automation.
A future Claude Code or competing local-AI instance that runs both on
cloud (1M token context for full-repository reasoning) and on local
hardware (small specialist agents for per-task prediction) would offer
both the power of cloud compute and the security and flexibility of
local computing for oncology trial sites and sponsors.

\clearpage

\subsection{Code Snippet: Hour-Loop Pattern (Simulation 3 Excerpt)}
\label{sec:methods-snippet}

The repeating pattern across the sponsor hourly scripts is a single
24-iteration loop that calls the 53 agent modules in dependency order
and writes the results to a per-hour JSON file. The following snippet
shows the simplified pattern that drives Simulation 3's 24-hour run.

\begin{lstlisting}
# Simulation 3 - sponsor/final_paper/scripts/hourly/sponsor_hour_00.py
# (simplified excerpt; full source has 53-agent dispatch + JSON output)
import json
from pathlib import Path
from core_agents import governance, study_execution, site_robotics, trust

OUTPUT_DIR = Path("output")
OUTPUT_DIR.mkdir(exist_ok=True)


def run_hour(hour_index: int) -> dict:
    state = {"hour": hour_index, "decisions": [], "escalations": []}
    state.update(governance.tick(hour_index))
    state.update(study_execution.tick(hour_index, state))
    state.update(site_robotics.tick(hour_index, state))
    state.update(trust.tick(hour_index, state))
    return state


if __name__ == "__main__":
    for hour in range(24):
        snapshot = run_hour(hour)
        out_path = OUTPUT_DIR / f"sponsor_hour_{hour:02d}.json"
        out_path.write_text(json.dumps(snapshot, indent=2))
\end{lstlisting}

The 168-hour Simulation 4 reuses the same hour-loop, with the loop
bound increased to 168 and the output directory renamed per day. The
local verification run on the Core i5-6200U laptop calls the same
\texttt{run\_hour} function with the same agent dispatch order, which
is why the smaller local output reflects the cloud output.
